\section{Methods}
\subsection{Linear Programming Framework}
Linear programming (LP) provides a mathematical framework for solving
optimization and feasibility problems subject to linear constraints. In the
present work, LP is used to reformulate activation screening as a feasibility
problem in composition space. Rather than generating candidate alloy compositions
and subsequently evaluating whether they satisfy prescribed activation criteria,
the activation criteria are expressed directly as constraints on alloy
composition. The set of compositions satisfying all constraints defines a
feasible region within the atomic-fraction simplex. This reformulation shifts the
focus from evaluating individual compositions to constructing and exploring the
admissible region of composition space itself.

The applicability of LP to the present problem arises from the linear dependence
of FLARE activation responses on alloy composition. As a result, the activation
criteria can be expressed as a system of linear inequalities whose solution set
forms a convex polytope. The geometry of this feasible space can be characterized
analytically through vertex enumeration, while feasible alloy compositions can be
obtained through LP-guided exploration of the polytope. The framework therefore
replaces conventional generate-and-filter screening with a constraint-first
methodology that directly identifies compositions satisfying the prescribed
activation requirements.
\subsubsection{Alloy design domain and FLARE response model}
The alloy design domain considered in this work consists of the eleven-element
refractory system comprising Ti, V, Ta, Nb, Mo, Zr, Cr, Hf, Fe, Re, and W. An
alloy is represented by its atomic-fraction vector
\(x=(x_1,\ldots,x_N)\) with \(N=11\), where \(x_i\) is the atomic fraction of
element \(i\). Physically admissible compositions lie on the standard
\((N-1)\)-simplex,

\[
\Delta=\Big\{x\in\mathbb{R}^{N}:\sum_i x_i=1,\; x_i\ge0\Big\}, \tag{1}
\]

and the discrete design domain is the regular 5 at.% grid on \(\Delta\), i.e.,
the set of compositions whose fractions are integer multiples of 0.05. On an
eleven-element simplex, this grid contains approximately \(3\times10^7\)
compositions, illustrating the computational burden associated with conventional
generate-and-filter screening.

Candidate alloys are screened on their neutron-activation response using the
FLARE reduced-order model. FLARE computes each activation response of an alloy as
a composition-weighted sum of pre-tabulated pure-element contributions, so that
every response is linear in alloy composition. Responses are evaluated at a
neutron flux of \(8.0\times10^{13}\), an irradiation time of 2 years, and the
decay times associated with each screening criterion.

Screening is performed using five activation responses: gas production
\(M_{\mathrm{GAS}}\), total heat output \(HT_T\), and displacement damage
\(DPA\) at a decay time of approximately 1 s, together with contact dose rate
\(DOSE\) at decay times of 1 month and 100 years. Each response combines from its
per-element contributions in the appropriate extensive basis -- by weight
fraction for the mass-extensive responses (\(M_{\mathrm{GAS}}\), \(HT_T\), and
\(DOSE\)) and by atomic fraction for the atom-extensive response (\(DPA\)). This
distinction determines the exact form of the linear constraints constructed in
Section 2.1.2. Activation thresholds are defined relative to pure W and EUROFER97
reference response vectors, as described in the LP formulation of the following
section.
\subsubsection{LP constraint construction}
\subsubsection{Feasible polytope and vertex enumeration}
\subsubsection{LP-guided composition enumeration}




